\section{From a local harmonic value to a theta rule}
\label{sec:local-theta-jets}

Let \(T\) be the complete discrete valuation ring of an integral chart,
\(K\) its fraction field and \(\pi\) a uniformizer. Write
\(v_T(\pi)=1\), \(v_T(0)=+\infty\), and
\(v_e(u)=v_T(u)/e\) for a positive integer \(e\). The chart carries
a truncated frame \(y\), the canonical connection, and a polarized
PEL uniformization with the specified isomorphism on torsion.
Write \(M\) for the positive level and \(\mathfrak m_b\) for its
chosen level ideal. The level condition includes
\(\mathfrak m_b+\overline{\mathfrak m}_b=\mathcal O_E\), where \(E\)
is the CM host. On the period lattice \(X\) take the integral
quadratic form \(Q\) furnished by \(\pi\), with
\(B(u,z)=Q(u+z)-Q(u)-Q(z)\).
Assume a basis \(b_j\) satisfies
\[
  Q(b_j)\le\Phi,\qquad |B(b_j,z)|\le M_0 Q(z)
       \quad(z\in X),\qquad \Phi\ge0 .
\]
On a zero lattice one can replace \(M_0\) by \(\max(M_0,0)\).
Let \(P\) be a harmonic polynomial of row degrees \(a_{\tau,i}\)
in \(n\ge1\) columns. Write \(\delta=\sum_{\tau,i}a_{\tau,i}\).
Every monomial in its support has degree \(\delta\).

For an integral power series \(f\) on the \(n\) copies of the chart,
write \(\mathcal D_{P,k}f\in K\) for its constant-term covariant
differential with polynomial \(P\) and line-bundle power \(k\).
It is linear in \(f\) and factors through the quotient by the
\((\delta+1)\)-st power of the ideal generated by the chart
coordinates. The coefficient field of \(P\) maps to \(K\) by the
fixed structural embedding. The local quantity occurring in the
cusp-jet estimate is
\[
       \nu_{\rm loc}(C;e,M,P)=
          \inf_{f\in T[[\boldsymbol Z]]}v_e(\mathcal D_{P,M}f).
\]
For \(x=(x_1,\ldots,x_n)\) in the discrete level-label set
\(\mathrm{La}^n\), denote the unique corresponding theta rule
by \(F_{x,P}\). Its evaluation in the frame \(y\) is in \(K\).

\paragraph{Two local hypotheses.}
The comparison below applies in either of the following settings:
\begin{enumerate}
\item The frame is algebraic, \(q=M\), \(M\) is a unit of \(T\),
\(\alpha\overline\alpha=M\) for \(\alpha\in\mathcal O_E\), and
\(\mathfrak m_b=(\overline\alpha)\).
\item The residue field of \(T\) has characteristic zero,
\(q\) is a divisor of \(M\), \(\alpha\overline\alpha=q\), and
\(\mathfrak m_b\subseteq(\overline\alpha)\). This formulation admits
the integral frame on the horizontal valuation without an algebraicity
hypothesis on that frame. Every nonzero rational integer is a unit of
\(T\) in this case.
\end{enumerate}
In both cases impose \(4\delta+4\le q\) and \(\delta M_0<q\). Then
\begin{equation}\label{eq:local-theta-comp}
 \inf_{x\in\mathrm{La}^n}v_e\bigl(F_{x,P}(y)\bigr)
 \ \le\ \nu_{\rm loc}(C;e,M,P)+\frac{\delta^2\Phi}{qe}.
\end{equation}
Both orders here are nonarchimedean; the value \(+\infty\) records a
zero evaluation. Positive residue characteristic \(p\) is permitted in
the first case when \(p\nmid M\). The Hermitian norm and its integral at
a complex embedding enter the metric comparison of
Section~\ref{sec:cusp-jets}.
\footnote{\href{https://github.com/ChristianKudeba/alpoge-rh-verification/blob/alpoge-source-7cbf7a0c/src/Development/AbelianVarietyTheory/Proofs/DivisorProofs/IInfThRuleLeNuLocOf.lean}{Full-level theorem and algebraic frame};
\href{https://github.com/ChristianKudeba/alpoge-rh-verification/blob/alpoge-source-7cbf7a0c/src/Development/AbelianVarietyTheory/Proofs/DivisorProofs/IInfThetaRuleLeLocalValueCharZero.lean}{divisor-level theorem in residue characteristic zero}.}

\paragraph{Changing the power in residue characteristic zero.}
In an integral biextension gauge, write \(S\) for the sum, over the
\(n\) columns, of the scalar gauge series at the distinguished
endomorphism. Its constant coefficient is zero. The formula for
\(\mathcal D_{P,k}\) applies the fixed differential to
\(f\exp(kS)\). The exponential is defined over \(T\), since \(T\)
contains \(\mathbb Q\). Put \(f'=f\exp((M-q)S)\). Each coefficient
is a finite sum in \(T\), and \(f'\exp(qS)=f\exp(MS)\).
The gauge-independence identity for harmonic \(P\) carries this equality
of differential values back to the chosen integral chart. Thus
\(\mathcal D_{P,q}f'=\mathcal D_{P,M}f\). Existence of the integral
biextension and its normalized gauge is used here. At full level
one simply takes \(f'=f\).
\footnote{\href{https://github.com/ChristianKudeba/alpoge-rh-verification/blob/alpoge-source-7cbf7a0c/src/Development/AbelianVarietyTheory/Proofs/DivisorProofs/HarmonicValueIntegralGermFramePowerCharZero.lean}{Integral gauge, power change, and harmonic gauge independence}.}

\paragraph{The small characters and the quadratic cost.}
Use the product period lattice \(X^n\) and
\(Q^{(n)}(z)=\sum_i Q(z_i)\). Putting \(b_j\) in the \(i\)-th
position gives a basis with the same bound \(\Phi\). Its mixed term
with \(z\) is \(B(b_j,z_i)\), bounded in absolute value by
\(\max(M_0,0)Q^{(n)}(z)\). In the following calculation \(Q,B,b_j\)
refer to this product and \(M_0\) to the resulting nonnegative bound.
Take \(\chi=\sum_j m_j b_j\), \(m_j\ge0\), \(\sum_jm_j\le\delta\).
The discriminant inequality for the positive quadratic form gives
\(\lvert B(s,t)\rvert\le 2\sqrt{Q(s)Q(t)}\): first use positivity
on the rational line \(s+ut\) by clearing denominators, then use
continuity in \(u\). Consequently
\[
 \sqrt{Q(\chi)}\le\sum_j m_j\sqrt{Q(b_j)}\le\delta\sqrt{\Phi},
 \qquad |B(\chi,z)|\le\delta M_0 Q(z).
\]
For \(z\ne0\), \(Q(z)\ge1\), whence
\[
 Q(\chi+qz)-Q(\chi)=q^2Q(z)+qB(\chi,z)
       \ge q(q-\delta M_0)Q(z)>0.
\]
Thus \(\chi\) is the unique minimum in \(\chi+qX^n\). This
strictness selects a single coefficient after reduction modulo \(\pi\).

Here are the geometric steps that turn these characters into all
jets. Let \(J=T[\boldsymbol Z]/(\boldsymbol Z)^{\delta+1}\).
Start on the abelian quotient of the semiabelian fiber. Write
\(q=4\delta+r\), \(r\ge4\), and fix a character \(\chi\).
The source's twisted-section lemma supplies a section
\(\beta_0\) of the \(\chi\)-twisted power \(r\) whose value in the
closed identity fiber is a unit. In degree four the very ample
sections \(s_i\) give an affine chart \(s_0\ne0\) at that point;
the ratios \(t_i=s_i/s_0\) generate the coordinate algebra of the
chart and their images generate the abelian quotient's jet algebra.
The latter is generated as a \(T\)-module by products of at most
\(\delta\) ratios. Indeed, write each ratio as its augmentation plus
an element of the augmentation ideal. All products of \(\delta+1\)
elements of that ideal vanish. Moreover \(t_0=1\), so shorter
products can be padded to length \(\delta\).
The factors \(s_{i_1}\cdots s_{i_\delta}\beta_0\) have precisely
the required power \(q\) and twist \(\chi\); in the local frame
their jets are the corresponding products of ratios times the fixed
unit \(s_0^\delta\beta_0\). Every abelian quotient jet therefore
occurs as a coefficient section \(\beta\) at \(\chi\).

Extend \(\beta\) by the functional equation of the degeneration in
the coset \(\chi+qX^n\), setting the other cosets to zero.
The coefficient valuation change on translation is
\((Q(\chi+qz)-Q(\chi))/q\). Thus the initial coefficient is
integral and the others in its coset are divisible by \(\pi\):
the change is positive and the coefficients are defined over the
discretely valued base. The integral expansion theorem realizes
this family as a section. Modulo \(\pi\) its jet is the abelian
coefficient times the torus character for \(\chi\), up to a fixed
unit from the frame. To see why these products span, the special
fiber torus supplies augmentation parameters of the form
\(v(b_j)/v(b_j)(0)-1\), where \(v(b_j)\) denotes the character
function in a trivialization. Together with the jets pulled back
from the abelian quotient these generate the augmentation ideal.
This is the tangent-space assertion for the semiabelian extension:
a tangent vector killed by the quotient differential is a torus
tangent, and the basis characters detect it. Passing successively
through the powers of the nilpotent augmentation ideal proves
generation of the jet algebra. Multiplying character functions
corresponding to \(b_j\)'s adds their characters up to a unit
pulled back from the abelian quotient. Expanding monomials of
degree at most \(\delta\) now produces the products indexed by the
\(\chi\)'s above. The fixed common frame unit preserves their span.
Nakayama's lemma lifts spanning modulo \(\pi\) to spanning over \(T\),
since \(J\) is finite free. A finite collection of the resulting
sections \(\sigma_i\) suffices. Each has expansion \(\theta_i\)
supported in the coset of one \(\chi_i\), with
\(Q(\chi_i)\le\delta^2\Phi\).
\footnote{\href{https://github.com/ChristianKudeba/alpoge-rh-verification/blob/alpoge-source-7cbf7a0c/src/Development/AbelianVarietyTheory/Proofs/JetHarmonics/SmallCombinationStrictMinMod.lean}{Strict coset minimum};
\href{https://github.com/ChristianKudeba/alpoge-rh-verification/blob/alpoge-source-7cbf7a0c/src/Development/AbelianVarietyTheory/Proofs/JetHarmonics/MemBasicOpenOfIsUnit.lean}{twisted abelian jets from degree-four sections};
\href{https://github.com/ChristianKudeba/alpoge-rh-verification/blob/alpoge-source-7cbf7a0c/src/Development/AbelianVarietyTheory/Proofs/JetHarmonics/IntegralFamilyOfStrictMin.lean}{integral extension in a single coset};
\href{https://github.com/ChristianKudeba/alpoge-rh-verification/blob/alpoge-source-7cbf7a0c/src/Development/AbelianVarietyTheory/Proofs/JetHarmonics/SmallCharactersSpanJets.lean}{special-fiber character spanning and integral lifting}.}

\paragraph{A scalar paying for all the characters at once.}
Extend \(v_T\) to the algebraic closure \(\overline K\), normalized
by \(v(\pi)=1\). In a frame of the \(q\)-th power on the truncated
chart, write \(f'\bmod(\boldsymbol Z)^{\delta+1}=\sum_i c_i
\operatorname{jet}(\sigma_i)\) with \(c_i\in T\).
The valuation condition on an expansion at level \(q\) and lower
bound zero means that the coefficient at \(\gamma\) has section
valuation at least \(Q(\gamma)/q\). For \(\theta_i\) the same
inequality holds with the lower bound decreased by \(Q(\chi_i)/q\):
it holds at the integral coefficient \(\chi_i\), propagates to its
coset by the functional equation, and holds at every zero coefficient.
Put \(k=\lfloor\delta^2\Phi\rfloor\). As \(Q(\chi_i)\) is an integer,
\(Q(\chi_i)\le k\) for every \(i\). Choose \(\lambda\in\overline K\)
with \(\lambda^q=\pi^k\). Then
\[
     0\le v(\lambda)=k/q\le\delta^2\Phi/q .
\]
The section \(\sigma=\sum_i\lambda c_i\sigma_i\) has expansion
\(s=\sum_i\lambda c_i\theta_i\) with lower bound zero. The
coefficients \(c_i\) are integral, and multiplying by \(\lambda\)
compensates every loss \(Q(\chi_i)/q\).
In the indicated frame its germ \(g\) satisfies
\(g\equiv\lambda f'\pmod{(\boldsymbol Z)^{\delta+1}}\).
Only one scalar is needed, regardless of the number of sections.
\footnote{\href{https://github.com/ChristianKudeba/alpoge-rh-verification/blob/alpoge-source-7cbf7a0c/src/Development/AbelianVarietyTheory/Proofs/JetHarmonics/ExistsPowFrame.lean}{The uniformized jet construction and floor of the quadratic bound};
\href{https://github.com/ChristianKudeba/alpoge-rh-verification/blob/alpoge-source-7cbf7a0c/src/Development/AbelianVarietyTheory/Proofs/JetHarmonics/ExistsGenOfTop.lean}{valuation propagation from a coset representative}.}

\paragraph{Selecting a translate.}
The containment of the level ideal and its conjugate implies
\((\alpha)+(\overline\alpha)=\mathcal O_E\). The invertibility
condition used by the theta-group argument is \(q\in T^\times\).
Choose the product theta frame and its pullback by \(\alpha\).
The integral spanning theorem in Section~\ref{sec:theta-translate-span}
expresses \(s=\sum_j d_j t_j\) with \(v(d_j)\ge0\). Each \(t_j\)
is a normalized translate of the pullback, at a point whose
coordinates are killed by \(\overline\alpha\).
The expansion bijection gives the section represented by each
\(t_j\); let \(h_j\) be its germ. Uniqueness and linearity of
expansion and germs give \(g=\sum_j d_jh_j\).
Write \(a_0=\mathcal D_{P,M}f\) and suppose \(a_0\ne0\).
Since the differential uses only the \(\delta\)-jet, after base
change to \(\overline K\) we have
\[
     \lambda a_0=\mathcal D_{P,q}g
                 =\sum_j d_j\mathcal D_{P,q}h_j .
\]
Some \(j\) therefore satisfies
\(v(\mathcal D_{P,q}h_j)\le v(a_0)+v(\lambda)\): if all summands had
larger order, so would the sum. The selected evaluation is nonzero.
\footnote{\href{https://github.com/ChristianKudeba/alpoge-rh-verification/blob/alpoge-source-7cbf7a0c/src/Development/AbelianVarietyTheory/Proofs/JetHarmonics/ExistsTermLe.lean}{Integral span of product translates and valuation selection};
\href{https://github.com/ChristianKudeba/alpoge-rh-verification/blob/alpoge-source-7cbf7a0c/src/Development/AbelianVarietyTheory/Proofs/DivisorProofs/JetsOnCopiesAny.lean}{product version for a general integral frame}.}

\paragraph{Identifying the value of the selected translate.}
The product-translation formula factors the selected germ, to order
\(\delta\), into a scalar \(u_0\) of valuation zero times the product
of the one-column germs. Each normalized one-column identification
contributes a root of unity \(u_i\); the normalized lifts have
finite order. Thus \(u=u_0\prod_i u_i\) has valuation zero.
Write \(M=qk'\). For a translated point \(b\) killed by
\(\overline\alpha\), the pullback by \(k'\overline\alpha\) takes the
translated power-\(q\) section to the power-\(M^2\) section defining
the level evaluation, with label representing \(\alpha b\).
The powers agree since \(q\,(k')^2\alpha\overline\alpha=M^2\).
The normalized pullbacks respect the level torsion trivializations;
applying the construction in each column gives an \(x\in\mathrm{La}^n\).
In the covariant evaluation the normalized chain rule for \(\beta=k'\overline\alpha\)
introduces, in a monomial of the given row degrees, the factor
\(N_a(\beta)\). The level definition divides the power-\(M^2\)
evaluation by \(M^\delta\). Since \(\alpha\beta=M\),
\(N_a(\alpha)N_a(\beta)=M^\delta\). Solving for the original
translated value gives the common factor
\[
    N_a(\alpha)=\prod_{\tau}\iota_\tau(\alpha)^{\sum_i a_{\tau,i}} .
\]
Here \(\iota_\tau\) denotes the corresponding embedding into the
coefficient field, viewed in \(\overline K\). Consequently the
chart-independent evaluation of the rule satisfies
\[
     \mathcal D_{P,q}h_j=u\,N_a(\alpha)\,F_{x,P}(y).
\]
Each \(\iota_\tau(\alpha)\) is integral over \(\mathbb Z\); an
element integral over \(\mathbb Z\) has nonnegative \(v\), by
applying the ultrametric inequality to its monic equation.
It follows that
\[
 v_T(F_{x,P}(y))\le v(\mathcal D_{P,q}h_j)
             \le v_T(a_0)+\delta^2\Phi/q .
\]
When \(a_0=0\) the comparison with this individual \(f\) holds in
the extended order as its right hand side is \(+\infty\). Divide by
\(e>0\) and take the infimum over all integral \(f\), obtaining
\eqref{eq:local-theta-comp}. The use in
Section~\ref{sec:cusp-jets} has \(e=1\).
\footnote{\href{https://github.com/ChristianKudeba/alpoge-rh-verification/blob/alpoge-source-7cbf7a0c/src/Development/AbelianVarietyTheory/Proofs/DivisorProofs/HarmonicValueTranslateDivisorLevelIntegralChart.lean}{Normalized pullback, the character factor, and its integrality};
\href{https://github.com/ChristianKudeba/alpoge-rh-verification/blob/alpoge-source-7cbf7a0c/src/Development/AbelianVarietyTheory/Proofs/JetHarmonics/HarmonicValueTranslateEqUnitMulNormaliserThetaVa.lean}{full-level version of the translate identification}.}

The comparison uses the polarized uniformization and torsion
identification on the prescribed chart, the normalized theta-lift and unit-ball properties in
Section~\ref{sec:theta-translate-span}, and the chart-independence
and pullback identities for the canonical connection. The preceding
calculation gives the interpolation cost and its transfer to a rule;
a derivation of those geometric inputs from the definitions is still
required for a standalone account. A finite place at which \(p\mid M\)
requires its own local argument since the integral spanning step above
uses that \(q\) is a unit.
